原统一链路通过 19 项原生测试、7 组 CPU／测试器用例和 4 组 XPU 用例。此后，SparseGate 分别通过合成输入 CPU 程序、预训练权重派生输入 CPU 程序及三处理源共存程序的验证。共存程序先执行原有 GPU／NPU 工作负载，再发出门控命令；该结果证明新增电路可与原处理源在保留的链路上共存，不代表门控与 GPU／NPU 同时争用资源。

每次门控运行均依据独立参考结果，逐项核对全部返回的得分／索引记录及聚集输出的每一个字节。小规模用例依次覆盖 FULL、有效 REUSE、REINDEX、过期 epoch 拒绝和 FULL 恢复。过期命令不产生被接受的 DMA 传输。验证交叉核对原生 AXI 波形、UCIe Flit、后端事件和存储命令，并保留可执行文件与源代码的摘要。每条命令的计数器均与 RTL 忙区间独立比对，成功 DONE 还须晚于最后一次 DMA 写响应。另设适配层与 RTL 联合测试，验证不支持的独占访问被拒绝，以及普通访问能够恢复。

波形检查器还识别出对已记录轨迹施加的三种受控错误：周期计数加一、DMA 读计数加一，以及在最后一次 B 通道握手时提前置位 DONE。上述错误仅通过检查时的数据覆盖注入，不修改原始仿真文件。

\begin{figure*}[t]
\centering\includegraphics[width=.94\textwidth]{completion_wave.pdf}
\caption{H4/N16/K4 CPU 程序中，首次成功 FULL 命令完成前后的实际原生 RTL VCD 采样。电平在模型求值后记录，握手判定采用上升沿之前稳定的信号值。DONE 在最后一次 B 通道握手后的第三个 AXI 周期提交。RTL STATUS 响应与其随后沿返回协议链路的交付记录独立关联。该图为仿真波形，并非芯片实测波形。}
\label{fig:completion}
\end{figure*}
\begin{table}[t]
\centering\caption{预训练权重派生 H32/N64/K8 程序的实测 RTL 命令计数。REINDEX 包含 16 个候选。R／W 为 DMA 读／写数据拍数，在两种存储周期下保持一致；每个选中索引对应搬运 288 B KV 数据。}
\label{tab:modes}
\small
\begin{tabular}{lrrrrr}\toprule
模式 & 评分条目数 & R & W & $1\times$ 周期 & $4\times$ 周期\\\midrule
FULL & 64 & 360 & 80 & 55,155 & 67,969\\
REUSE & 0 & 72 & 80 & 1,949 & 5,783\\
REINDEX & 16 & 218 & 80 & 18,788 & 26,359\\
\bottomrule\end{tabular}
\end{table}
AXI 仿真周期为 2 ns。命令周期数计量 RTL 从开始忙碌到接收最后写响应并提交完成状态的区间，主机程序的运行时间则还包含协议传输与软件交互。上述输入下的三条成功命令均返回八条结果记录。

将在线存储周期由 0.25 ns 增大至 1 ns 后，被测可执行文件、RTL、得分和 DMA 数量均保持不变，所有成功命令的周期数均增加，完整主机程序结束时间增加 181.530 $\mu$s。轮询请求数量可以随等待时间变化，因此该实验用于确认在线存储延迟能够反馈到电路及主机执行过程，不能解释为固定指令工作量下的加速比比较。

另设完整配置 H32/N640/K512 程序，先后执行 FULL 和 REUSE，并使用彼此独立的目标地址。每条命令均核对全部 512 条得分／索引记录、填充字节及 147,456 B 聚集数据。FULL 需要 670,572 个忙周期，产生 6,624 个读数据拍和 5,120 个写数据拍；REUSE 需要 115,776 个周期，读／写数据拍数分别为 4,608／5,120。该用例将完整配置的选择验证扩展至实际在线链路；小规模 REINDEX 与错误恢复序列由另一组用例覆盖。

\begin{figure*}[t]
\centering\includegraphics[width=.86\textwidth]{native_modes.pdf}
\caption{两种在线存储周期下 H32/N64/K8 命令的实测延迟，以及两种周期下相同的 DMA 有效载荷。REUSE 保留已有选择，REINDEX 对外部提供的候选子集重新算分；两者工作内容不同，不能作为可互换的精度基线。有效载荷包含结果记录与聚集 KV 数据的写回，不含主机上传、MMIO、回读及链路封装开销。}
\label{fig:native}
\end{figure*}
